% 调整页边距 不要动！
\usepackage{fullpage}
\usepackage[top=0.7cm, bottom=4.5cm, left=1.5cm, right=1.5cm]{geometry}
% 调整页边距 不要动！
\usepackage[T1]{fontenc}
\usepackage[utf8]{inputenc}
\usepackage[english]{babel}
\usepackage{csquotes}
\usepackage[style=alphabetic,backend=biber]{biblatex}
\usepackage{arydshln}
\addbibresource{../.knowledge/references.bib}
% 数学和物理包 不要动！
\usepackage{amsmath,amsthm,amsfonts,amssymb,amscd}
\usepackage{physics}
% \usepackage{cite}
\usepackage{booktabs}
% 数学和物理包 不要动！
\usepackage{longtable}
\usepackage{geometry}
% 其他必须的包 不要动！
\usepackage{enumerate}
\usepackage{fancyhdr}
\usepackage{mathrsfs}
\usepackage{algorithm}
\usepackage{algpseudocode}
\usepackage{graphicx}
\graphicspath{{figures/}}
% 其他必须的包 不要动！
% 当前环境使用 pdfLaTeX 时，ctex 的中文字体集会因为缺少可用字体而中断编译。
% 本文不依赖中文正文，因此改为保留兼容命令并关闭中文字体加载。
\newcommand{\heiti}{\bfseries}
\usepackage{pythonhighlight} %Python代码环境包，使用时将Python代码粘贴到\begin{python}...\end{python}
\usepackage{indentfirst}
% 调整节标题 不要动
% \usepackage{titlesec}
% \renewcommand{\thesection}{\arabic{section}}
% \titleformat{\section}
% {\normalfont\bfseries}{\large Question~\thesection}{1em}{}
% 调整节标题 不要动

% 调整段间距 可以自行修改
\setlength{\parindent}{0em}
\setlength{\parskip}{0.25em}
% 调整段间距


% 调整页眉 可以自行修改
\newcommand\course{\textbf{Reports}}   % 页眉右侧内容
\newcommand\problemnumber{2}    % 第几题
\newcommand\probletitle{\textit{\heiti Fault-tolerant quantum computation using Rydberg atoms}}  %题目标题
\newcommand\timeID{\textbf{\today}}          
% 调整页眉

% 设置页眉格式 不要动！
\pagestyle{fancy}
\headheight 25pt
% \lhead{\includegraphics[scale = 0.1]{Figures/jlulogo.png}}  % 添加页眉logo
% \rhead{\course \\ \timeID}
\rhead{\course}
\lfoot{}
\cfoot{}
\rfoot{\small\thepage}
\headsep 1.5em
% 设置页眉格式 不要动！

% 设置段落缩进和行距
\linespread{1.1}

% 设置段落缩进和行距

% 设置图片注释字体大小
\usepackage[font=scriptsize]{caption}
% \usepackage[english, chinese]{babel}

\usepackage{hyperref}


% \newcommand{\bra}[1]{\langle #1\rvert}
% \newcommand{\ket}[1]{\lvert #1\rangle}
\newcommand{\proj}[1]{\lvert #1\rangle\!\langle #1\rvert}
% \newcommand{\braket}[2]{\langle #1\vert #2\rangle}
\newcommand{\SYC}[1]{\color{blue}(#1)\color{black}}

% General Commands
%%%%%%%%%%%%%%%%%%%%%%%%%%%%%%%%%%%%%%%%%%%%%%%%%%%%%%%%%%%%%%%%%%%%%%%%%%%%%%%%%%%%%%%%%%%%%%%%%%%%%%%%%%%%%%%%%%%%%%%
\newtheorem{theorem}{Theorem}[section]
\newtheorem{proposition}[theorem]{Proposition}
\newtheorem{property}[theorem]{Property}
\newtheorem{definition}{Definition}[section]
\newtheorem{lemma}{Lemma}
\newtheorem{remark}[theorem]{Remark}
\newtheorem{corollary}[theorem]{Corollary}
\newtheorem{observation}[theorem]{Observation}
\newtheorem{example}[theorem]{Example}
\newtheorem{exercise}{Exercise}
\newtheorem{problem}{Problem}
